\documentclass[a3paper, landscape]{article}
\usepackage[left=2cm, right=2cm, top=2cm, bottom=2cm]{geometry}
\usepackage[T2A]{fontenc}
\usepackage[utf8]{inputenc}
\usepackage[russian]{babel}
\thispagestyle{empty}

% The tikz-er2 package.
% Ver. 1.0
% (c) Pavel Caldo
% January 2009

\NeedsTeXFormat{LaTeX2e}
\ProvidesPackage{tikz-er2}[2009/01/31 Entity-relationship diagrams using the
Tikz library]

\RequirePackage{tikz}

\usetikzlibrary{shapes.geometric}
\usetikzlibrary{arrows}
\usetikzlibrary{fit}

\tikzstyle{every entity} = []
\tikzstyle{every weak entity} = []
\tikzstyle{every attribute} = []
\tikzstyle{every relationship} = []
\tikzstyle{every link} = []
\tikzstyle{every isa} = []

\tikzstyle{link} = [>=triangle 60, draw, thick, every link]

\tikzstyle{total} = [link, double, double distance=3pt]

\tikzstyle{entity} = [rectangle, draw, black, very thick,
                      minimum width=6em, minimum height=3em,
                      every entity]

\tikzstyle{weak entity} = [entity, double, double distance=2pt,
                           every weak entity]

\tikzstyle{ident relationship} = [relationship, double, double distance=2pt]

\tikzstyle{attribute} = [ellipse, draw, thin, black,
                         minimum width=4em, minimum height=2em,
                         text width=4em, align=center,
                         every attribute]

\tikzstyle{multi attribute} = [attribute, double, double distance=1.5pt]

\tikzstyle{derived attribute} = [attribute, dashed]

\tikzstyle{key attribute} = [attribute]

\tikzstyle{relationship} = [diamond, draw, black, very thick,
                            minimum width=6em, minimum height=4em,
                            aspect=1.5,
                            every relationship]

\tikzstyle{isa} = [regular polygon, regular polygon sides=3,
                   draw, black, very thick,
                   minimum width=4em, minimum height=4em,
                   every isa]

\newcommand{\key}[1]{\underline{#1}}

\usetikzlibrary{positioning}

\tikzstyle{every entity} = [fill=white]
\tikzstyle{every attribute} = [fill=white, node distance=10em, text width=6.5em, minimum width=6em]
\tikzstyle{every relationship} = [fill=white, aspect=1.5]
\tikzstyle{every edge} = [link]

\begin{document}
\begin{figure}[p]
\centering
\scalebox{0.80}{
\begin{tikzpicture}[node distance=6em, every edge/.style={link}]

% ИЗДАТЕЛЬСТВО
\node[entity] (izd) {Издательство};
\node[attribute] (i1) [above left of=izd] {\key{Название}} edge (izd);
\node[attribute] (i2) [above of=izd] {Страна} edge (izd);
\node[attribute] (i3) [above right of=izd] {Сайт} edge (izd);
\node[attribute] (i4) [left of=izd] {Год основания} edge (izd);
\node[attribute] (i5) [below left of=izd] {Email} edge (izd);

% выпускает
\node[relationship] (r_vyp) [right=8em of izd] {Выпускает}
    edge node[above] {1} (izd);

% ГАРНИТУРА
\node[entity] (garn) [right=8em of r_vyp] {Гарнитура};
\draw[link] (r_vyp) -- node[above] {N} (garn);
\node[attribute] (g1) [above left of=garn] {\key{Название}} edge (garn);
\node[attribute] (g2) [above of=garn] {Год создания} edge (garn);
\node[attribute] (g3) [above right of=garn] {Описание} edge (garn);
\node[attribute] (g4) [below right of=garn] {Число начертаний} edge (garn);

% относится к
\node[relationship] (r_otn) [right=8em of garn] {Классифицируется}
    edge node[above] {N} (garn);

% распространяется под
\node[relationship] (r_ras) [below left=12em of garn] {Распространяется}
    edge node[above] {1} (garn);

% КАТЕГОРИЯ
\node[entity] (kat) [right=8em of r_otn] {Категория};
\draw[link] (r_otn) -- node[above] {1} (kat);
\node[attribute] (k1) [above left of=kat] {\key{Название}} edge (kat);
\node[attribute] (k2) [above of=kat] {Описание} edge (kat);
\node[attribute] (k3) [above right of=kat] {Иконка} edge (kat);
\node[derived attribute] (k4) [right of=kat] {Число гарнитур} edge (kat);
% выдаёт
\node[relationship] (r_vyd) [below=8em of izd] {Выдаёт}
    edge node[left] {1} (izd);

% ЛИЦЕНЗИЯ
\node[entity] (lic) [below=8em of r_vyd] {Лицензия};
\draw[link] (r_vyd) -- node[left] {N} (lic);
\draw[link] (r_ras) -- node[left] {N} (lic);
\node[attribute] (l1) [above left of=lic] {\key{Тип}} edge (lic);
\node[attribute] (l2) [left of=lic] {Срок действия} edge (lic);
\node[attribute] (l3) [below left of=lic] {Макс. станций} edge (lic);
\node[attribute] (l4) [below of=lic] {Стоимость} edge (lic);
\node[attribute] (l5) [below right of=lic] {Дата приобр.} edge (lic);

% включает
\node[relationship] (r_vkl) [below=8em of garn] {Включает}
    edge node[left] {1} (garn);

% НАЧЕРТАНИЕ
\node[entity] (nach) [below=8em of r_vkl] {Начертание};
\draw[link] (r_vkl) -- node[left] {N} (nach);
\node[attribute] (n1) [above right of=nach] {\key{Название}} edge (nach);
\node[attribute] (n2) [below right=2em and 8em of nach] {Насыщенность} edge (nach);
\node[attribute] (n3) [below left of=nach] {Наклон} edge (nach);
\node[attribute] (n4) [left of=nach] {Имя файла} edge (nach);
\node[attribute] (n5) [above left of=nach] {Размер
файла} edge (nach);

% хранится в
\node[relationship] (r_hr) [right=14em of nach] {Хранится}
    edge node[above] {N} (nach);

% ФОРМАТ
\node[entity] (fmt) [right=14em of r_hr] {Формат};
\draw[link] (r_hr) -- node[above] {1} (fmt);
\node[attribute] (f1) [above of=fmt] {\key{Название}} edge (fmt);
\node[attribute] (f2) [above right of=fmt] {Расширение} edge (fmt);
\node[attribute] (f3) [right of=fmt] {Год появления} edge (fmt);
\node[attribute] (f4) [below right of=fmt] {Описание} edge (fmt);
\node[attribute] (f5) [below of=fmt] {Поддержка веба} edge (fmt);

% поддерживает
\node[relationship] (obl) [below=12em of nach] {Содержит}
    edge node[left] {1} (nach);

% ПОДДЕРЖКА ЯЗЫКА
\node[entity] (podder) [below=8em of obl] {Алфавит};
\draw[link] (obl) -- node[left] {N} (podder);
\node[derived attribute] (pl2) [above left of=podder] {Число начертаний в системе} edge (podder);
\node[derived attribute] (pl3) [below left of=podder] {Число поддерж. гарнитур} edge (podder);
\node[attribute] (pl4) [below of=podder] {Дата проверки} edge (podder);
\node[attribute] (pl5) [below right of=podder] {Статус} edge (podder);

% поддерживает
\node[relationship] (r_pod1) [left=12em of podder] {Охватывает}
    edge node[left] {N} (podder);


% ЯЗЫК
\node[entity] (yaz) [below=8em of r_pod1] {Язык};
\draw[link] (r_pod1) -- node[left] {1} (yaz);
\node[attribute] (y1) [above left of=yaz] {\key{Название}} edge (yaz);
\node[attribute] (y2) [left of=yaz] {Код ISO} edge (yaz);
\node[attribute] (y3) [below left of=yaz] {Направление письма} edge (yaz);
\node[attribute] (y4) [below of=yaz] {Тип письменности} edge (yaz);


% ИСПОЛЬЗОВАНИЕ (entity)
\node[entity] (isp) [below right=18em and 18em of nach] {Установка};
\node[attribute] (u1) [above of=isp] {Дата установки} edge (isp);
\node[attribute] (u2) [below left of=isp] {Дата удаления} edge (isp);
\node[attribute] (u3) [left of=isp] {Путь к файлу} edge (isp);
\node[attribute] (u4) [below right of=isp] {Статус} edge (isp);
\node[attribute] (u5) [below of=isp] {Дата посл. исп.} edge (isp);

% Начертание -- фиксирует -- Использование
\node[relationship] (r_fix) [below right=8em and 8em of nach] {Участвует};
\draw[link] (nach) -- node[above left] {1} (r_fix);
\draw[link] (r_fix) -- node[above right] {N} (isp);

% размещено на
\node[relationship] (r_razm) [right=8em of isp] {Производится}
    edge node[above] {N} (isp);

% РАБОЧАЯ СТАНЦИЯ
\node[entity] (rab) [right=8em of r_razm] {Рабочая станция};
\draw[link] (r_razm) -- node[above] {1} (rab);
\node[attribute] (r1) [above left of=rab] {\key{Инв. номер}} edge (rab);
\node[attribute] (r2) [above of=rab] {Процессор} edge (rab);
\node[attribute] (r3) [above right of=rab] {Объём памяти} edge (rab);
\node[attribute] (r4) [right of=rab] {Тип накопителя} edge (rab);
\node[attribute] (r5) [below right of=rab] {Графический адаптер} edge (rab);

% установлена
\node[relationship] (r_ust) [below=8em of rab] {Установлена}
    edge node[left] {N} (rab);

% ОПЕРАЦИОННАЯ СИСТЕМА
\node[entity] (os) [below=4em of r_ust] {ОС};
\draw[link] (r_ust) -- node[left] {1} (os);
\node[attribute] (o1) [left of=os] {\key{Название}} edge (os);
\node[attribute] (o2) [below left of=os] {Версия} edge (os);
\node[attribute] (o3) [below of=os] {Разрядность} edge (os);
\node[attribute] (o4) [below right of=os] {Год выпуска} edge (os);

\node[text width=20em, anchor=north west, font=\small] (goals) at ([xshift=-15em, yshift=0]current bounding box.north east) {
    \textbf{Цели функционирования базы данных:}
    \begin{enumerate}
        \item Хранение сведений о шрифтовых гарнитурах, их начертаниях и характеристиках.
        \item Хранение информации о рабочих станциях и их аппаратно-программных конфигурациях.
        \item Хранение информации о лицензиях на шрифты, условиях использования и сроках действия.
        \item Хранение информации об установленных шрифтах на каждой рабочей станции.
    \end{enumerate}

    \textbf{Текст чтения диаграммы:}
    \begin{enumerate}
        \item Издательство выпускает гарнитуру.
        \item Издательство выдаёт лицензию.
        \item Гарнитура классифицируется по категории.
        \item Гарнитура включает начертание.
        \item Начертание сожерит алфавит.
        \item Алфавит охватывает язык.
        \item Начертание участвует в установке.
        \item Установка производится на рабочей станции.
        \item На рабочую станцию установлена операционная система.
        \item Начертание хранится в формате.
        \item Гарнитура распространяется по лицензии.
    \end{enumerate}
};

\end{tikzpicture}
}
\caption{ER-диаграмма предметной области <<Учёт компьютерных шрифтов в электронной типографии>>}
\end{figure}
\end{document}
